% appendix_priority.tex: priority of the theoretical results (2026-09-09).
% Companion: memos/priority_ledger_2026-09-09.md (table, verification status, missing bib entries).
% Citations without a key in references.bib are written in plain text; the entries are listed in the memo.
% Input in main.tex after appendix_strategies.

\section{Priority of the Theoretical Results}\label{appendix:priority}

This appendix records, for each theoretical result in the paper, the closest prior statement of the insight and the increment the paper makes. The purpose is to place the results rather than to claim them: several are special cases of known principles, and the appendix says which.

\paragraph{Sensitivities to the payment schedule (Proposition \ref{prop:sensT}).} The device of interpreting a ratio of two behavioral responses as a ratio of marginal utilities is that of \citet{chetty_moral_2008}, who separates the liquidity and moral-hazard components of the response of job search to a benefit by comparing it with the response to a cash grant at the same date, and of \citet{indarte_moral_2023}, who applies the same comparison to bankruptcy filing with a dollar of payment relief and a dollar of seizable equity, both valued at the filing date (her equation (14)). \citet{ganong_liquidity_2020} supply the two-date empirical comparison, payment reduction against principal reduction, and interpret it through liquidity and wealth. In each case the transfers compared are located at one date or are not indexed by their date. Proposition \ref{prop:sensT} locates the transfer at an arbitrary future date in the optimal-stopping default problem and shows that the response factors into the discount function $D(s-1)$, the probability $Q_s$ of still being in good standing, and the average marginal utility of a repayer at $s$ relative to that of the marginal defaulter today; the two-period identity $A/(A+B) = \delta\lambda$ of Section \ref{sec:suffstats} is its $T=2$ case. \citet{campbell_model_2015} compute the same comparative statics numerically in a lifecycle model; the closed-form factorization is the paper's.

\paragraph{The present-payment premium (Proposition \ref{prop:premium}).} That a liquidity-constrained household values a dollar today above a dollar later by a ratio of marginal utilities, and that a comparison of a cash transfer with a later or an illiquid transfer measures that ratio, is the content of the liquidity-constraint literature from \citet{zeldes_consumption_1989} to \citet{chetty_moral_2008} and \citet{indarte_moral_2023}, whose ratio of responses is a ratio of marginal utilities at the filing date. Proposition \ref{prop:premium} places that ratio in the default first-order condition: the payment due now is priced at the marginal defaulter's marginal utility and every later payment at the repayers' average, so the ratio $\lambda_s$ multiplies the discount function in the weight on each later payment and is bounded away from one for a constrained marginal defaulter at any discount rate. The increment is the location of the ratio in the weights, which is what makes it inseparable from present bias by payment timing (Theorem \ref{thm:betalambda}).

\paragraph{Identification of the discount curve (Proposition \ref{prop:curveT}).} The recovery of a discount function at several horizons from dated choices is the program of the elicitation literature: \citet{coller_eliciting_1999}, \citet{harrison_estimating_2002} and \citet{andreoni_estimating_2012} elicit it from monetary choices, \citet{benhabib_present-bias_2010} test its shape across horizons, \citet{warner_personal_2001} recover one rate from a large-stakes lump-sum-or-annuity choice, and \citet{laibson_estimating_2024} estimate $\beta$ and $\delta$ by matching lifecycle consumption and credit moments in a structural model. None uses default responses, and all either hold survival at one or model it inside the structure. Proposition \ref{prop:curveT} states that within one population the ratio of default responses to two payment-path perturbations is the ratio of their weighted sums with the common scale cancelled, so that with survival measured from repayment data and the marginal-utility ratio supplied the discount function is identified pointwise and the exponential null is testable with three or more horizons. The two-arm designs of \citet{ganong_liquidity_2020} and \citet{dobbie_targeted_2020} are treated in Section \ref{sec:quasi} as pairs of points on that curve, an interpretation neither paper makes; the two-period statistic of the earlier literature is one point on it.

\paragraph{Interior default and the degeneracy of committed-repayment comparisons (Proposition \ref{prop:interiorT}, Lemma \ref{lem:committedT}).} Default as an option exercised each period against an exclusion or flow penalty, with interior default probabilities generated by income shocks, is the model of \citet{eaton_debt_1981}, of \citet{chatterjee_quantitative_2007}, who prove existence with interior default, and, in a finite-horizon lifecycle setting, of \citet{campbell_model_2015}. The backward unraveling of repayment when the sanction vanishes in the last period is the argument of \citet{bulow_sovereign_1989}. Proposition \ref{prop:interiorT} is a special case of these existence results, stated for a fixed payment schedule with the recursion $K_t = \sigma + \delta\,\mathbb{E}[(K_{t+1} - g_{t+1}(y))^+]$, which gives interior probabilities for every horizon and every schedule under a flow penalty and $u(0^+) = -\infty$. The increment is Lemma \ref{lem:committedT}: replacing the option to reconsider with a commitment to all remaining installments drives the default probability to one at long horizons, so that the corner probabilities of the committed-repayment model are a consequence of that modeling choice, and a penalty proportional to the remaining balance inherits the endgame problem that a flow component removes.

\paragraph{Non-separability of present bias and the marginal-utility ratio (Theorem \ref{thm:betalambda}, Corollary \ref{cor:separate}; Appendix \ref{appendix:nonclassical}).} That utility curvature and liquidity confound money-based measures of discounting and present bias is established by \citet{andersen_eliciting_2008}, who require a joint risk elicitation to correct for curvature, by \citet{andreoni_estimating_2012}, who identify curvature and discounting jointly from convex budgets, and by \citet{augenblick_working_2015}, who show that present bias measured in money is contaminated by liquidity and curvature and measure it in effort instead; \citet{frederick_time_2002} and \citet{cohen_measuring_2020} catalogue the confounds. \citet{laibson_estimating_2024} separate $\beta$ from $\delta$ structurally from the joint holding of liquid debt and illiquid wealth, and \citet{kuchler_sticking_2021} identify present bias from paydown plans against their realization. The paper's statement is the exact form of the confound in the default setting: under $D(k) = \beta\delta^k$ and a constant $\lambda$, every ratio $(\partial p_1/\partial m_s)/(\partial p_1/\partial m_1)$ equals $\beta\lambda\,\delta^{s-1} Q_s$, so $\beta$ and $\lambda$ enter only as a product and no perturbation of a current payment against later ones separates them. The constructive half is the observation that the pair of responses in \citet{indarte_moral_2023} compares currencies at a fixed date rather than dates, and therefore identifies $\lambda$ of equation \eqref{eq:twoperiod}; the structural separation in \citet{laibson_estimating_2024} rests on the same kind of comparison.

\paragraph{Beliefs and preferences in the weights (Proposition \ref{prop:beliefs}).} That the effective discount factor on date-$s$ utility is the product of the pure discount function and the subjective probability of reaching $s$ is \citet{yaari_uncertain_1965}; that choice data alone cannot separate preferences from expectations, so that expectations must be measured, is \citet{manski_measuring_2004}. The paper applies both to the default option, where survival is survival in good standing through $s$ and is endogenous to the default, payoff and cure hazards, states that the discount function is recovered only with the borrower's subjective survival set equal to the measured $Q_s$, and measures the inputs to $Q_s$ from the credit-bureau panel: the auto exit hazard and the auto and mortgage cure rates. In the population of \citet{ganong_liquidity_2020}, measured survival rather than time preference accounts for most of the apparent discount in the constant-weight estimate (Section \ref{sec:quasi}).

\paragraph{The term margin as two channels and exposure (Proposition \ref{prop:termmargin}; Section \ref{sec:termmargin}, Appendix \ref{appendix:balancepenalty}).} Maturity in consumer credit has been studied as a demand object and as a screening object. \citet{attanasio_credit_2008} and \citet{karlan_credit_2008} show that constrained borrowers' demand responds to maturity far more than to the interest rate; \citet{hertzberg_screening_2018} show that maturity choice screens on private default risk; \citet{argyle_monthly_2020} show that borrowers target the monthly payment and choose maturity jointly with it; \citet{adams_liquidity_2009} and \citet{einav_contract_2012} show that loan size raises default and is chosen under private information; \citet{ghent_recourse_2011} show that recourse lowers mortgage default among underwater borrowers. None decomposes the default response to maturity at a fixed payment. The paper's contribution is the decomposition of that response into a preference channel discounted at the borrower's rate, a balance-scaled penalty channel discounted at the contract rate, and, for lifetime outcomes, an exposure term; the prediction that the penalty channel dominates at short horizons and reverses as the loan amortizes, which produces the horizon profile and the score gradient of Section \ref{sec:termmargin}; and the use of state deficiency-judgment thresholds on auto loans, alongside the recourse states of \citet{ghent_recourse_2011} for mortgages, as a test of the penalty channel in installment credit.

\paragraph{Exact aggregation (Proposition \ref{prop:aggT}).} That a pooled regression or instrumental-variables coefficient under heterogeneity is a weighted average of local responses with explicit non-negative weights is established by \citet{yitzhaki_using_1996}, \citet{angrist_two-stage_1995} and \citet{heckman_structural_2005}; that sufficient-statistic formulas aggregate heterogeneous responses with the welfare-relevant weights is the point of \citet{chetty_sufficient_2009} and of \citet{kleven_sufficient_2021}. Proposition \ref{prop:aggT} is the closed form for the ratio statistic of this model: the pooled estimate is the average of type-level statistics with weights equal to each type's population share times its density at its own default margin times the responsiveness of its marginal defaulter, so that it lies in the convex hull of the type statistics and is the discount factor of the population at the default margin. The form is elementary once the two sensitivity formulas are in hand; its use is the interpretation of the level of the payment elasticity as exposure in Section \ref{sec:heterogeneity}.

\paragraph{The run-off of the payment margin (Remark \ref{rem:option}, Appendix \ref{appendix:mcliquidity}).} That the loan-age profile of default hazards is governed by option values is the finding of \citet{deng_mortgage_2000} and of \citet{campbell_model_2015}; that continuation values embed the option, so that responses to future payoffs are muted relative to their discounted value, is the logic of the dynamic discrete-choice literature following \citet{rust_optimal_1987}; \citet{fuster_payment_2017} measure the default response to payment size at a fixed loan age. The paper's result is negative, established by simulation rather than theorem, and conditional on the calibration: in the $T$-period model calibrated to the sample's default rate the ratio of the payment elasticity in the last loan-age bin to the first lies between $0.77$ and $0.82$ at every discount rate on a grid from $0$ to $1.2$ per year, and between $0.20$ and $0.22$ with a persistent income process, in both cases at a payment elasticity ten times the data's (Appendix \ref{appendix:persistent}); within those calibrations the shape of the run-off contains no information about the rate, and a first-order formula that includes survival and the marginal-utility ratio but not the option value attributes the option's truncation of far payments to impatience. Structural default papers, including \citet{chatterjee_quantitative_2007} and \citet{athreya_persistence_2019}, report the sensitivity of the level of default to the discount factor; none, to the author's knowledge, reports the insensitivity of an elasticity profile to it.
